			\section{计算几何}
				\subsection{二维几何基础操作}
				%\subsection{凸包}
					\inputminted{cpp}{src/Geometry/凸包.cpp}
					\inputminted{cpp}{src/Geometry/闵可夫斯基和.cpp}
					\inputminted{cpp}{src/Geometry/geo.cpp}
				%	\inputminted{cpp}{src/Geometry/far_pair.cpp}
				% \subsection{直线半平面交}
					\inputminted{cpp}{src/Geometry/半平面交.cpp}
				\subsection{整数半平面交}
					\inputminted{cpp}{src/Geometry/integral_hpi.cpp}
				\subsection{凸包询问：凸包内、切点、交点、最近点}
					\inputminted{cpp}{src/Geometry/convex_findmax.cpp}



				 \subsection{点、线段在简单多边形内}
					\inputminted{cpp}{src/Geometry/AirportConstruction.cpp}
				\subsection{$O(n ^ 2)$ Ear Clipping 三角剖分}
				\inputminted{cpp}{src/Geometry/Triangulation.cpp}
				\subsection{单插入动态凸包}
				\inputminted{cpp}{src/Geometry/dynamic_convex_hull.cpp}
					% \subsection{凸包快速询问}
					% \inputminted{cpp}{src/Geometry/logconvexhull.cpp}
%				\subsection{直线与凸包交点}
%				\inputminted{cpp}{src/Geometry/直线与凸包交点.cpp}
%				\subsection{点到凸包切线}
%				\inputminted{cpp}{src/Geometry/点到凸包切线.cpp}
				% \subsection{Delaunay 三角剖分}
				% \inputminted{cpp}{src/Geometry/DelaunayTriangulation.cpp}

				\subsection{圆}
					\inputminted[highlightlines={7}]{cpp}{src/Geometry/circle.cpp}
				\subsection{圆反演，阿波罗尼茨圆}
				\begin{small}
					\input{src/Geometry/阿波罗尼茨圆.tex}
				\end{small}
					\inputminted{cpp}{src/yzh/circle_inversion.cpp}
				\subsection{圆并}
					\inputminted{cpp}{src/Geometry/圆并.cpp}
				\subsection{多边形与圆交}
					\inputminted{cpp}{src/Geometry/多边形和圆的交.cpp}

				%\subsection{圆幂~圆反演~根轴}
				%	\input{src/Geometry/圆幂.tex}
				\subsection{球面基础, 经纬度球面距离}
				\begin{small}
					\input{src/Geometry/球面.tex}
				\end{small}
					\inputminted{cpp}{src/Geometry/经纬度求球面最短距离.cpp}
%				\subsection{最近点对}
%				\inputminted{cpp}{src/Geometry/最近点对.cpp}
				%\subsection{长方体表面两点最短距离}
				%\inputminted{cpp}{src/Geometry/长方体表面两点最短距离.cpp}
				\subsection{圆上整点}
					\inputminted{cpp}{src/Geometry/圆上整点.cpp}
				\subsection{三角形}
				\inputminted{cpp}{src/Geometry/三角形.cpp}
				\subsection{相关公式：三角形心，海伦公式，欧拉定理，内接球，外接圆}
					\input{src/Geometry/Formula(Saga).tex}
				\subsection{三维几何基础操作}
					\inputminted{cpp}{src/Geometry/三维几何.cpp}
				\subsection{三维凸包}
					\inputminted{cpp}{src/Geometry/三维凸包.cpp}
				\subsection{最小覆盖球}
					\inputminted{cpp}{src/Geometry/最小覆盖球.cpp}
				%	\newpage
\subsection{黄金三分}
					\inputminted[highlightlines={7}]{cpp}{src/yzh/golden_ternary.cpp}
\subsection{自适应 Simpson}
					\inputminted{cpp}{src/Math/Simpson.cpp}
